\documentclass[12pt]{article}
\usepackage[utf8]{inputenc}
\usepackage[T1]{fontenc}
\usepackage{hyperref}
\title{Bindal Language Resources for K-12 Curriculum Development}
\author{Compiled by Claude Sonnet 5.5 (claude-sonnet-5-5) for Andr\'as Kornai}
\date{}
\begin{document}
\maketitle

\section{Introduction}
Bindal (xbd) is an extinct Aboriginal language of North Queensland. Its country ran from Cape Cleveland toward Ayr and the mouth of the Burdekin River, within today's Townsville and Burdekin council areas. Only a few word lists survive, and classification is uncertain: Bowern suggests it may be Maric, while Breen considers it one of the Lower Burdekin languages. Wikipedia gives a few words, as spelt by Bindal authors, for example \emph{gamu} ``water'', \emph{barri} ``stone'' and \emph{adha} ``yes''. Neighbours include the Juru and the Gia. I found no language programme, so a curriculum can only be a cultural and heritage unit with a few recovered words.

\section{Primary Educational Resources}
None found.

\section{Linguistic and Reference Materials}
\url{https://en.wikipedia.org/wiki/Bindal_language}; \url{https://en.wikipedia.org/wiki/Bindal_people}; \url{https://en.wikipedia.org/wiki/Lower_Burdekin_languages}. The South Australian Museum archives list a Bindal language group page (contents not checked): \url{https://samuseum.sa.gov.au/collection/archives/language_groups/bindal}.

\section{Religious and Community Resources}
The Bindal and Wulgurukaba peoples are named as Traditional Owners of the Townsville area in the results (native title claims in 2016, and a 2022 reconciliation steering committee at a local quarry). I found no Bindal-specific community language organisation. Neighbouring peoples: \url{https://en.wikipedia.org/wiki/Juru_people} and \url{https://Www.Wikipedia.org/wiki/Gia_people}. A regional site, ``What Mob Is That?'', also appeared: \url{https://www.aboutthenorth.au/styled-59/styled-131/styled-184/styled-238/styled-655/}.

\section{Digital and Audio Resources}
None found.

\section{Curriculum Development Context}
Any revitalization work would need to be led by Bindal Traditional Owners, using archival wordlists and comparison with related Maric languages. Local Townsville and Burdekin schools could use a place-name and country-focused unit.
\end{document}